% !TeX root = ./../main.tex

\chapter{Wprowadzenie}\label{chapter:introduction}

\agentinfo{
Rozdział przedstawia przyczyny powstania pracy, problem badawczy, tezy, cele, wkład oraz układ rozprawy.
Ma charakter opisowy.
Szczegółowe uzasadnienie luki badawczej znajduje się w rozdziale przeglądu literatury \agentref{chapters/literature_review.tex}, a tutaj należy się do niego odwołać.
Rozdział pisać po rozdziałach 2--9, a wcześniej utrzymywać go jako aktualny konspekt.
Docelowa objętość: około 6 stron; cała praca około 125 stron.

Podstawa treści (ustalenie autora, 2026-11-26): wstępy dostarczonych artykułów \agentref{artykuły/Quantum_Errors_ICCS/main.tex} \agentref{artykuły/PROCS_KES2023_LATEX_Template/k23-558.tex} \agentref{artykuły/CoDIT_2025_final_paper/root.tex} \agentref{artykuły/Quantum_journal_GAN__Copy_/article.tex} oraz wprowadzenie rozprawy Michała jako wzór zakresu i toku wywodu \agentref{doktorat Michał/chapters/introduction.tex}.
Z rozprawy Michała nie przepisywać zdań.
Jego praca dotyczy kwantowego kodowania obrazów i opisuje ten temat szeroko; tutaj kodowanie obrazów jest tylko tłem, a ciężar leży na błędach i ich ograniczaniu.
Nowe pozycje literatury tylko prawdziwe, sprawdzone w źródle przed zacytowaniem.

Oś narracji całej pracy:
(1) komputery ery NISQ są zaszumione, a pełna korekcja błędów zużywa kubity i bramki, których brakuje;
(2) kwantowe przetwarzanie obrazów wymaga wielu kubitów i głębokich obwodów, więc jest szczególnie wrażliwe na szum;
(3) autor proponuje ograniczanie błędu po stronie klasycznej, na podstawie pomiarów: metodę PDU;
(4) PDU zostaje przeniesiona na kodowanie i rekonstrukcję obrazów metodą LPIQE oraz porównana z FRQI i NEQR;
(5) błąd pozostały po PDU jest dalej redukowany siecią generatywną pix2pix;
(6) ta sama idea korekty opartej na pomiarze jest sprawdzana dla uczonego kodowania obrazów na współczesnym urządzeniu;
(7) badany jest wpływ szumu i redukcji błędów na klasyfikację;
(8) perspektywa uzupełniająca, jako prace powiązane z udziałem autora: szum ma strukturę, którą można modelować (generatory łamiące symetrię) i wykorzystać w uczeniu (sieci oparte na grupach Liego).
}

\section{Motywacja}\label{sec:introduction:motivation}
\agentinfo{
Treść: ograniczenia sprzętu ery NISQ (liczba kubitów, objętość kwantowa, dekoherencja), koszt kwantowej korekcji błędów, potencjał kwantowego przetwarzania obrazów (wykładnicza pojemność rejestru), napięcie między liczbą kubitów przeznaczonych na obliczenia a przeznaczonych na korekcję.
Kontekst aplikacyjny z artykułów: platformy fotoniczne, komunikacja kwantowa, rozpoznawanie przyrządów lotniczych w symulatorze lotu.
Źródła: wstępy artykułów \agentref{artykuły/Quantum_Errors_ICCS/main.tex} \agentref{artykuły/PROCS_KES2023_LATEX_Template/k23-558.tex} \agentref{artykuły/CoDIT_2025_final_paper/root.tex} \agentref{artykuły/Quantum_journal_GAN__Copy_/article.tex} oraz koncepcja rozprawy \agentref{inne/Wstępna koncepcja rozprawy doktorskiej.docx}.
Dane liczbowe o sprzęcie z koncepcji z 2022 r. są nieaktualne i wymagają aktualizacji na podstawie bieżących źródeł.
}

\section{Problem badawczy, tezy i cele pracy}\label{sec:introduction:theses}
\agentinfo{
Tezy z koncepcji rozprawy z 2022 r. (punkt wyjścia):
T1. Metody redukcji błędów kwantowych oparte na pomiarze, w tym funkcja PDU, w sposób statystycznie istotny zmniejszają wpływ dekoherencji i szumu powstających w obliczeniach kwantowych.
T2. Błąd średniokwadratowy obrazu zrekonstruowanego po zastosowaniu PDU w kwantowym przetwarzaniu obrazów jest na poziomie standardowych metod kompresji obrazów.

Robocza propozycja rozszerzenia, DO UZGODNIENIA Z PROMOTOREM (nie wpisywać do tekstu bez akceptacji autora):
T3. Połączenie PDU z klasycznym modelem generatywnym typu obraz--obraz pozwala dalej zmniejszyć błąd rekonstrukcji obrazu bez użycia dodatkowych zasobów kwantowych.
T4. Klasyczna redukcja błędów oparta na pomiarze poprawia skuteczność klasyfikacji prowadzonej na wynikach zaszumionych obliczeń kwantowych; teza do sprawdzenia w nowych eksperymentach na symulatorze, wynik może być negatywny.
Wcześniejsza robocza teza o strukturze szumu odpada: prace z S. Raubitzkiem są pracami powiązanymi z udziałem autora, a nie wkładem rozprawy (ustalenie autora z 2026-11-26).
Autor dopuszcza dorobienie badań, preferencyjnie na symulatorze, aby tezy miały pokrycie w wynikach.

Robocze pytania badawcze:
PB1. Czy błąd obliczeń opartych na próbkowaniu kwantowym jest stabilny między kalibracjami urządzenia w stopniu pozwalającym wyznaczyć jego korektę eksperymentalnie?
PB2. Czy korekta wyznaczona dla funkcji referencyjnej poprawia wyniki innych obliczeń na tym samym układzie kubitów?
PB3. Jak duża jest utrata danych w metodach FRQI, NEQR i LPIQE na symulatorze i na rzeczywistych urządzeniach oraz w jakim stopniu ogranicza ją PDU?
PB4. Czy klasyczna sieć generatywna może usunąć błąd pozostały po PDU i jakim kosztem?
PB5. Jak redukcja błędów wpływa na skuteczność klasyfikacji obrazów?
PB6. Jak uczone kodowanie obrazów znosi szum rzeczywistego urządzenia i czy korekta kalibracyjna zmniejsza lukę między symulatorem a sprzętem?

DECYZJE AUTORA z 2026-11-27 (po przeglądzie literatury i wynikach rozdziałów 7 i 8):
PDU jest metodą rozwijaną w ramach tego doktoratu i ma być opisana jako ważny wkład rozprawy.
Jej wyróżnikami względem metod pokrewnych z lat 2020--2023 (klucze jattana_general_2020, czarnik_error_2021, urbanek_mitigating_2021, lolur_reference-state_2023) są: korekta indeksowana znanym argumentem, wyznaczana raz dla rodziny obwodów, oraz zbadanie stabilności błędu między kalibracjami urządzenia; metody pokrewne trzeba przy tym przywołać i nie wolno pisać, że ich nie ma (zestawienie w sekcji 3.3, tabela tab:lr:qem:calibrated).
Robocza teza T4 odpada: klasyfikacja zostaje pytaniem badawczym (PB5) z odpowiedzią warunkową z rozdziału 8, bez osobnej tezy.
W luce badawczej elementy (a)--(d) mają węższe brzmienie, a element o strukturze szumu usunięto (konspekt sekcji sec:lr:gap).
Uwaga orkiestratora do rozważenia przez autora (to nie jest decyzja): teza T1 w brzmieniu ogólnym obejmuje wszystkie metody redukcji oparte na pomiarze, a wyniki rozdziałów 7 i 8 dla korekty typu PDU poza obliczeniami QCoSamp są w większości negatywne, więc T1 warto zawęzić do zakresu, w którym PDU zbadano w rozdziale 4.

Teza T2 wymaga doprecyzowania: z jaką metodą kompresji i przy jakich parametrach następuje porównanie.
Tytuł pracy obejmuje klasyfikację obrazów, a dotychczasowe artykuły jej nie pokrywają; pokrycie dają nowe eksperymenty opisane w rozdziale 8 \agentref{chapters/classification.tex}.
}

\section{Zakres pracy}\label{sec:introduction:scope}
\agentinfo{
Wskazać, co praca obejmuje: obliczenia oparte na próbkowaniu kwantowym, redukcję błędów po stronie klasycznej, obrazy monochromatyczne małych rozmiarów (8x8 na urządzeniach, 16x16 na symulatorze), urządzenia nadprzewodzące IBM Quantum oraz symulatory.
Wskazać, czego praca nie obejmuje: projektowania nowych kodów korekcyjnych, eksperymentów na sprzęcie fotonicznym (LPIQE jest projektowana dla tej platformy, ale badania wykonano na urządzeniach nadprzewodzących i symulatorach).
Wyjaśnić użycie terminów ,,ograniczanie'' (redukcja), ,,mitygacja'' i ,,korekcja'' błędów oraz odesłać do definicji w rozdziale 2 \agentref{chapters/theoretical_background_03_QEC.tex}.
}

\section{Wkład pracy}\label{sec:introduction:contribution}
\agentinfo{
Wkład należy wyliczyć w punktach i powiązać z rozdziałami:
(a) metoda PDU: definicja, procedura kalibracji, dowód eksperymentalny stabilności i istotności statystycznej (rozdział 4);
(b) uogólnienie PDU na przetwarzanie obrazów i schemat kodowanie--rekonstrukcja z LPIQE (rozdział 5);
(c) porównanie utraty danych w FRQI, NEQR i LPIQE na symulatorze i rzeczywistych urządzeniach (rozdział 5);
(d) zastosowanie sieci pix2pix do redukcji błędu pozostałego po PDU (rozdział 6);
(e) charakterystyka szumu i korekta kalibracyjna dla uczonego kodowania obrazów na rzeczywistym urządzeniu (rozdział 7);
(f) ocena wpływu szumu i redukcji błędów na klasyfikację, w tym klasyfikację obrazów (rozdział 8).
Prace omawiane w rozdziale 9 nie są wkładem rozprawy; autor implementował w nich część kodu i uruchamiał eksperymenty.
W punktach (e) i (f) oddzielić wkład własny autora od wkładu współautorów manuskryptów roboczych.
}

\section{Publikacje i manuskrypty związane z rozprawą}\label{sec:introduction:publications}
\agentinfo{
Status publikacji (opublikowana, przyjęta, w recenzji, manuskrypt) oraz pełne dane bibliograficzne wymagają potwierdzenia przez autora.
Stan zweryfikowany w Crossref 2026-11-26: opublikowane z autorem jako pierwszym autorem są trzy prace (ICCS 2022, KES 2023, CoDIT 2025); praca o sieciach GAN istnieje tylko jako preprint z 2023 r.; praca o szumie, w której autor jest drugim autorem, została opublikowana w APS Open Science w 2026 r.; praca o grupach Liego jest preprintem z 2026 r.
Lista jest zgodna z profilem autora w serwisie ResearchGate odczytanym 2026-11-26 (10 pozycji, w tym preprinty i praca magisterska).
Praca FRQI Pairs (CoDIT 2025, autor rozprawy jako czwarty autor) dotyczy klasyfikacji obrazów; ustalić z autorem jego udział i to, czy warto ją omówić w rozdziale 8 lub w przeglądzie literatury.
Rozprawa powstaje z wyników eksperymentów, a nie z artykułów, więc lista poniżej dokumentuje dorobek, a nie wyznacza zakres pracy.
Dwa kolejne manuskrypty robocze autora: klasyfikator z odczytem fazowym (repozytorium quantum-classifier) oraz hybrydowy autoenkoder obrazów z uczonym kodowaniem (katalog best_results_package, plik article/main.tex); oba współautorskie i nieopublikowane.
Wpisy znajdują się w bibliography/own_publications.bib; pola oznaczone DO WERYFIKACJI uzupełnić po potwierdzeniu.
Manuskrypt o sieciach GAN istnieje w dwóch szablonach (Quantum oraz MDPI Energies); ustalić, gdzie ostatecznie został zgłoszony.
Artykuł CoDIT 2025 ma dwie wcześniejsze wersje robocze w katalogach Quantum_Errors_ICCS_Images i Artykuł_2025_PDU_conf; wersją obowiązującą jest CoDIT_2025_final_paper.
}

\subsection{Prace bezpośrednio związane z rozprawą}

Autor rozprawy jest pierwszym autorem wszystkich prac wymienionych w tym punkcie.

\begin{itemize}
\item Praca~\cite{werner_experiment_driven_2022} wprowadza metodę redukcji błędów opartą na funkcji PDU i stanowi podstawę rozdziału~\ref{chapter:pdu}.
\item Praca~\cite{werner_data_loss_2023} porównuje utratę danych w metodach kwantowej reprezentacji obrazów i stanowi podstawę rozdziału~\ref{chapter:pdu_images}.
\item Praca~\cite{werner_error_reduction_2025} przedstawia schemat kodowania i~rekonstrukcji obrazów z~redukcją błędów i stanowi podstawę rozdziału~\ref{chapter:pdu_images}.
\item Preprint~\cite{werner_deep_learning_manuscript} opisuje zastosowanie generatywnych sieci przeciwstawnych do redukcji błędów i stanowi podstawę rozdziału~\ref{chapter:dl_error_reduction}.
\end{itemize}

\subsection{Pozostałe prace}

\begin{itemize}
\item Praca~\cite{raubitzek_quantum_noise} dotyczy rekonstrukcji szumu kwantowego z~generatorów łamiących symetrię i jest omawiana w rozdziale~\ref{chapter:noise_structure}.
\item Preprint~\cite{raubitzek_lie_group_manuscript} dotyczy sieci neuronowych opartych na grupach Liego i jest omawiany w rozdziale~\ref{chapter:noise_structure}.
\item Praca~\cite{potempa_frqi_pairs_2025} dotyczy klasyfikacji obrazów kodowanych metodą FRQI z~użyciem kwantowej sieci rekurencyjnej.
\item Praca magisterska autora~\cite{werner_magisterka_2022} zawiera pierwotne sformułowanie metody PDU.
\end{itemize}

\section{Struktura pracy}\label{sec:introduction:structure}
\agentinfo{
Przewodnik po rozdziałach: po jednym akapicie na rozdział, z odwołaniami przez etykiety chapter:... .
Pisać na końcu.
}
